\chapter{Optymalizacja i testy}
\label{chap:optymalizacja}

Końcowym etapem było przeprowadzenie audytu konfiguracji, czyszczenie śladów oraz wyeksportowanie maszyny do obrazu \texttt{VulnBox.ova} wraz z testem out-of-the-box.

\section{Zautomatyzowany audyt konfiguracji i usług}
\label{sec:verify-vm}

Aby wyeliminować ryzyko pominięcia któregokolwiek z komponentów przed eksportem obrazu, opracowano skrypt \texttt{box/scripts/verify\_vm.sh}. Skrypt ten realizuje kontrolę plików konfiguracyjnych, uprawnień i gniazd sieciowych (\texttt{ss -tln}) oraz lokalne testy ścieżek wobec usług na portach \texttt{22/tcp}, \texttt{445/tcp}, \texttt{8080/tcp} oraz \texttt{5000/tcp}, w tym wywołanie \texttt{?proxy\_path=} i weryfikację nadpisania modułu przez \texttt{PYTHONPATH} z poziomu konta \texttt{dev}.

\begin{lstlisting}[language=bash, caption={Wynik działania skryptu audytowego verify\_vm.sh na maszynie VulnBox}, label={lst:verify-output}]
$ sudo bash /tmp/APNWA-VulnBox/box/scripts/verify_vm.sh
=== VulnBox End-to-End Verification Suite ===
[PASS] Hostname is set to 'vulnbox'
[PASS] System account 'dev' exists
[PASS] Netplan uses 01-netcfg.yaml (mode 0600) with no installer/cloud-init overrides
[PASS] GRUB2 custom PBKDF2 configuration present in /etc/grub.d/40_custom
[PASS] Active /boot/grub/grub.cfg contains password_pbkdf2 directive
[PASS] Active /boot/grub/grub.cfg enables --unrestricted unattended boot
[PASS] UFW firewall is active
[PASS] UFW rule present for 22/tcp
[PASS] UFW rule present for 445/tcp
[PASS] UFW rule present for 8080/tcp
[PASS] OpenSSH listener (ssh.socket / ssh.service) is active
[PASS] Systemd service 'fail2ban' is active
[PASS] Systemd service 'smbd' is active
[PASS] Systemd service 'nginx' is active
[PASS] Systemd service 'flask-internal' is active
[PASS] Systemd service 'cron' is active
[PASS] SSH directory permissions and sshd PubkeyAuthentication readiness verified
[PASS] Listening sockets verified for :445 (Samba) and 127.0.0.1:5000 (Flask API)
[PASS] Samba configuration valid and deployment-notes.txt present
[PASS] PHP-FPM socket /run/php/php-fpm.sock is present
[PASS] HTTP probe on 127.0.0.1:8080/?module=status returned portal status
[PASS] Stage 2 path normalization check retrieved /opt/internal_api/config.ini
[PASS] Internal Flask API /health endpoint responded on 127.0.0.1:5000
[PASS] Internal Flask API /api/v1/info authenticated using token from config.ini
[PASS] Web proxy bridge 127.0.0.1:8080/?proxy_path= relayed request to Flask API
[PASS] Sudoers policy /etc/sudoers.d/90-dev-admin passed visudo syntax check
[PASS] Cron policy (0644 root:root), backup_check.sh (+x), and pre_backup.sh (0775 root:dev) valid
[PASS] Module /opt/admin/status_helper.py permissions are 0664 (root:dev)
[PASS] Sudoers PYTHONPATH override loaded custom status_helper.py as user dev
[PASS] User flag /home/dev/user.txt exists with mode 0600 and owner dev
[PASS] Root flag /root/root.txt exists with mode 0600 and owner root

=== Summary: 31 passed, 0 failed ===
\end{lstlisting}

Uzyskanie wyniku \texttt{31 passed, 0 failed} (listing \ref{lst:verify-output}) potwierdza pełną spójność środowiska przed przystąpieniem do procedury czyszczenia.

\section{Higiena obrazu przed eksportem}
\label{sec:cleanup-vm}

Po zakończeniu audytu \texttt{verify\_vm.sh} oraz ewentualnych testów manualnych z rozdziału \ref{chap:weryfikacja} konieczne jest przywrócenie stanu początkowego wszystkich modyfikowanych plików tak, aby osoba importująca obraz \texttt{VulnBox.ova} otrzymała czyste środowisko bez pozostawionych kluczy SSH, logów testowych czy historii poleceń. Zadanie to automatyzuje skrypt \texttt{box/scripts/cleanup\_vm.sh} (uruchamiany obowiązkowo po \texttt{verify\_vm.sh}, bezpośrednio przed \texttt{sudo poweroff}), który wykonuje osiem operacji:
\begin{enumerate}
    \item zatrzymanie usług generujących logi (\texttt{cron}, \texttt{fail2ban}, \texttt{nginx}, \texttt{smbd}, \texttt{flask-internal}) na czas czyszczenia,
    \item wyzerowanie pliku \path{/home/dev/.ssh/authorized_keys} i przywrócenie uprawnień \texttt{0600} (\texttt{dev:dev}),
    \item odtworzenie oryginalnych wersji plików \path{/opt/admin/hooks/pre_backup.sh} (\texttt{0775}, \texttt{root:dev}) oraz \path{/opt/admin/status_helper.py} (\texttt{0664}, \texttt{root:dev}) wraz z usunięciem katalogów \texttt{\_\_pycache\_\_},
    \item bezpieczne wyzerowanie aktywnych plików logów usług (\texttt{nginx}, \texttt{fail2ban}, \texttt{auth.log}, \texttt{syslog}, \texttt{dpkg}, \texttt{wtmp}/\texttt{btmp}/\texttt{lastlog}) oraz usunięcie plików w \path{/var/log/samba/}, \path{/var/log/apt/} i archiwów rotacyjnych,
    \item usunięcie nieużywanych pakietów, pamięci podręcznej i list APT (\texttt{apt-get autoremove --purge}, \texttt{apt-get clean}, \path{/var/lib/apt/lists/*}) oraz zredukowanie dziennika \texttt{journalctl --vacuum-time=1s},
    \item usunięcie plików historii powłoki i sesji (\texttt{.bash\_history}, \texttt{.lesshst}, \texttt{.viminfo}) dla konta \texttt{root} oraz wszystkich katalogów domowych w \path{/home/*},
    \item wyczyszczenie katalogów tymczasowych \path{/tmp} i \path{/var/tmp},
    \item wyzerowanie wolnego miejsca na partycji głównej (\texttt{dd if=/dev/zero}) z zachowaniem bufora \texttt{1024 MB} oraz uruchomienie \texttt{fstrim} przed wyłączeniem maszyny w celu zmniejszenia rozmiaru eksportowanego pliku \texttt{VulnBox.ova}.
\end{enumerate}

\section{Eksport do pliku VulnBox.ova i test out-of-the-box}
\label{sec:eksport-ova}

Po wyłączeniu maszyny poleceniem \texttt{sudo poweroff} odłączono pomocniczy interfejs NAT oraz wirtualny napęd ISO, pozostawiając wyłącznie kartę sieci laboratoryjnej. Następnie wyeksportowano maszynę do standardu OVF 2.0 \cite{ovf20} przy użyciu narzędzia \texttt{VBoxManage}:

\begin{lstlisting}[language=bash, caption={Polecenie eksportu maszyny wirtualnej do pliku VulnBox.ova}, label={lst:vbox-export}]
$ VBoxManage export "VulnBox" --output "VulnBox.ova" --ovf20 --vsys 0 \
    --product "VulnBox Educational Appliance" \
    --description "APNWA Project: Vulnerable Virtual Machine (Ubuntu Server 26.04.1 LTS)"
\end{lstlisting}

Wygenerowany plik \texttt{VulnBox.ova} zaimportowano testowo pod nową nazwą w programie z zaznaczoną opcją reinicjalizacji adresów MAC wszystkich kart sieciowych. Podczas rozruchu potwierdzono:
\begin{itemize}
    \item brak możliwości wejścia w tryb edycji parametrów jądra w menu GRUB2 bez podania hasła administratora, przy jednoczesnym automatycznym starcie systemu bez ingerencji użytkownika,
    \item poprawne przypisanie nowego adresu IPv4 z serwera DHCP dzięki uniwersalnej regule \texttt{match: name: "en*"} w konfiguracji Netplan,
    \item automatyczny start wszystkich sześciu kluczowych usług (\texttt{ssh}, \texttt{fail2ban}, \texttt{smbd}, \texttt{nginx}, \texttt{flask-internal}, \texttt{cron}) natychmiast po załadowaniu systemu.
\end{itemize}
